\chapter{Binary Search and Unimodal Optimisation}\label{ch:binsearch}

Binary search finds a boundary in a monotone world in logarithmically many steps. Everyone knows it as the way to look up a
value in a sorted array, but its real power in contests lies elsewhere: in \emph{binary search on the answer}, which turns
``find the best value'' into the far easier question ``is this value achievable?'', and in its relatives for real numbers,
unimodal functions, ratios and batches of queries. This chapter develops binary search from a single loop invariant, so
that the off-by-one errors that plague it disappear, and then applies it in all of these forms.

\begin{chapterintro}{\bronze\ \silver\ \gold\ \ioi}{Chapters~\ref{ch:complexity} and~\ref{ch:sorting}.}
  \item write binary search from an invariant, and use \texttt{lower\_bound} and \texttt{upper\_bound} correctly;
  \item recognise monotone predicates and binary search on the answer;
  \item find the $k$-th smallest element of a set that is never built;
  \item binary search over real numbers, and maximise ratios with fractional programming;
  \item optimise unimodal and convex functions with ternary search;
  \item answer many binary-search queries at once with parallel binary search.
\end{chapterintro}

% ======================================================================================================================
\section[Searching a sorted array]{Searching a sorted array \hfill \bronze}\label{sec:bsearch-array}

To find where a value $x$ belongs in a sorted array, compare it with the middle element and discard the half that cannot
contain the answer. Each comparison halves the range, so after $\lceil \log_2(n + 1)\rceil$ comparisons nothing is left:
20 comparisons suffice for $10^6$ elements and 60 for $10^{18}$.

The cleanest way to write it, and to never get it wrong, is to decide exactly what the answer is and to maintain an
\emph{invariant}. We look for the first index $i$ with $a_i \ge x$ (or $n$ if there is none), and keep a half-open range
$[lo, hi)$ with the invariant
\[
  a_i < x \text{ for all } i < lo, \qquad a_i \ge x \text{ for all } i \ge hi .
\]
Initially $lo = 0$ and $hi = n$, and the invariant holds vacuously. Look at $mid = \lfloor (lo + hi)/2 \rfloor$, which lies in
$[lo, hi)$: if $a_{mid} < x$, the invariant still holds with $lo = mid + 1$; otherwise it holds with $hi = mid$. Either way the
range shrinks, and when $lo = hi$ the invariant says that $lo$ is exactly the first index with $a_i \ge x$
(Figure~\ref{fig:bsearch-steps}).

\begin{figure}[htbp]
\centering
\begin{tikzpicture}
  \foreach \row/\lo/\hi/\mid/\txt in {0/0/10/5/{$a_5 = 23 \ge 23$: \ $hi \gets 5$},
      1/0/5/2/{$a_2 = 8 < 23$: \ $lo \gets 3$}, 2/3/5/4/{$a_4 = 16 < 23$: \ $lo \gets 5$}} {
    \pgfmathsetmacro{\yy}{-1.25*\row}
    \foreach \v [count=\i from 0] in {2, 5, 8, 12, 16, 23, 38, 56, 72, 91} {
      \pgfmathtruncatemacro{\inside}{(\i >= \lo && \i < \hi) ? 1 : 0}
      \ifnum\i=\mid \node[cellc, hot] at (\i*0.7, \yy) {\v};
      \else\ifnum\inside=1 \node[cellc] at (\i*0.7, \yy) {\v};
      \else \node[cellc, dim] at (\i*0.7, \yy) {\v}; \fi\fi
    }
    \node[blbl, anchor=east] at (-0.55, \yy) {step \the\numexpr\row+1\relax};
    \node[notel, anchor=west] at (7.2, \yy) {\txt};
    \draw[line width=1.2pt, draw=fgblue] ({\lo*0.7 - 0.35}, \yy + 0.45) -- ({\lo*0.7 - 0.35}, \yy - 0.45);
    \draw[line width=1.2pt, draw=fgblue] ({\hi*0.7 - 0.35}, \yy + 0.45) -- ({\hi*0.7 - 0.35}, \yy - 0.45);
  }
  \foreach \i in {0,...,9} \node[lbl] at (\i*0.7, -3.15) {\i};
  \node[lbl, anchor=east] at (-0.55, -3.15) {index};
  \node[notel, anchor=west] at (7.2, -3.15) {$lo = hi = 5$: the first $a_i \ge 23$};
\end{tikzpicture}
\caption{Binary search for the first element $\ge 23$. The blue bars delimit the range $[lo, hi)$ that may still contain the
boundary; grey cells are already decided, and red marks the middle element compared in each step.}
\label{fig:bsearch-steps}
\end{figure}

\subsection{The standard library}

The library functions implement exactly this search, and it is worth knowing precisely what each returns
(Figure~\ref{fig:lowerupper}):
\begin{itemize}
  \item \cc{lower_bound(all(a), x)}: the first position with value $\ge x$, the insertion point of $x$ before any equal
        values;
  \item \cc{upper_bound(all(a), x)}: the first position with value $> x$, the insertion point after equal values;
  \item \cc{equal_range(all(a), x)}: both at once; the difference of the two positions is the number of copies of $x$;
  \item the last value $\le x$ (the \emph{predecessor}) is just before \cc{upper_bound}, if that is not the beginning, and
        the first value $\ge x$ (the \emph{successor}) is at \cc{lower_bound}, if that is not the end.
\end{itemize}
All of them return iterators; subtract \cc{a.begin()} to get an index. A reversed array is searched with
\cc{greater<>()} as the comparator, and \cc{binary_search} only answers yes or no, which is rarely enough.

\begin{figure}[htbp]
\centering
\begin{tikzpicture}
  \arrs{A}{0}{0}{1/n, 3/cool, 3/cool, 3/cool, 5/n, 8/n, 8/n, 9/n}
  \idx{A}{8}
  \draw[->, line width=1.1pt, draw=fgblue] ($(A-2.north)+(0,1.25)$) node[above, font=\sffamily\scriptsize\bfseries,
    text=fgblue, inner sep=1.5pt] {lower\_bound(3) = 1} -- ($(A-2.north)+(0,0.04)$);
  \draw[->, line width=1.1pt, draw=fgverm] ($(A-5.north)+(0,0.55)$) node[above, font=\sffamily\scriptsize\bfseries,
    text=fgverm, inner sep=1.5pt] {upper\_bound(3) = 4} -- ($(A-5.north)+(0,0.04)$);
  \node[lbl] at ($(A-3.south)+(0,-0.7)$) {count of 3 = $4 - 1 = 3$};
  \node[notel, anchor=west, text width=5.8cm] at (6.1, 0) {For $x = 4$, which does not occur, both functions return
    4: the position where 4 would be inserted. The predecessor of 4 is the element before it (3), and its successor is the
    element at that position (5).};
\end{tikzpicture}
\caption{\texttt{lower\_bound} and \texttt{upper\_bound} on a sorted array with repeated values. The copies of $x$ occupy
exactly the positions from \texttt{lower\_bound} up to (not including) \texttt{upper\_bound}.}
\label{fig:lowerupper}
\end{figure}

\begin{problem}{Counting values in ranges}
Given $n$ numbers and $q$ queries $[l, r]$, report how many of the numbers lie in each range. Limits: $n, q \le 2\cdot10^5$.
\solution Sort once. The values in $[l, r]$ occupy the positions from the first value $\ge l$ up to (not including) the first
value $> r$, so each query is two binary searches: $\Oh{(n + q)\log n}$ in total.
\end{problem}

\begin{code}
void solve() {
    int n, q;
    cin >> n >> q;
    vi a(n);
    cin >> a;
    sort(all(a));
    while (q--) {
        ll l, r;
        cin >> l >> r;
        // [first value >= l, first value > r) holds the values in [l, r]
        auto from = lower_bound(all(a), l), to = upper_bound(all(a), r);
        print(to - from);
    }
}
\end{code}

\begin{pitfall}[Searching a \texttt{set} with the free function]
\cc{lower_bound(all(s), x)} compiles for a \cc{set} or \cc{map}, but set iterators cannot jump to the middle, so the free
function walks through the elements one by one and takes \emph{linear} time. Always use the member function
\cc{s.lower_bound(x)}, which descends the tree in $\Oh{\log n}$. This single mistake turns many correct solutions into time
limit verdicts.
\end{pitfall}

% ======================================================================================================================
\section[Binary search on a monotone predicate]{Binary search on a monotone predicate \hfill \silver}\label{sec:predicate}

Nothing in the argument above used the array itself, only the fact that the condition $a_i \ge x$ is false for a prefix of
the indices and true for the rest. So binary search works for any \defn{monotone predicate}: a yes/no function
$\mathit{ok}(v)$ on integers such that once it is true, it stays true for all larger $v$ (Figure~\ref{fig:monotone}). The
search finds the boundary, the first $v$ with $\mathit{ok}(v)$ true, with $\Oh{\log(\text{range})}$ evaluations of
$\mathit{ok}$, and the predicate can be anything you can compute: a greedy simulation, a flow, a DP.

\begin{figure}[htbp]
\centering
\begin{tikzpicture}
  \foreach \x in {0,...,11} {
    \pgfmathtruncatemacro{\lab}{\x + 10}
    \ifnum\x>4 \node[cellc, good] at (\x*0.72, 0) {T}; \else \node[cellc, dim] at (\x*0.72, 0) {F}; \fi
    \node[lbl] at (\x*0.72, -0.55) {\lab};
  }
  \node[blbl, anchor=east] at (-0.55, 0) {$\mathit{ok}(v)$};
  \node[lbl, anchor=east] at (-0.55, -0.55) {$v$};
  \draw[->, line width=1.1pt, draw=fgverm] (5*0.72, 1.05) node[above, font=\sffamily\scriptsize\bfseries, text=fgverm]
    {first true: the answer} -- (5*0.72, 0.4);
  \begin{scope}[yshift=-1.9cm]
    \foreach \x in {0,...,11} {
      \ifnum\x<8 \node[cellc, good] at (\x*0.72, 0) {T}; \else \node[cellc, dim] at (\x*0.72, 0) {F}; \fi
    }
    \node[blbl, anchor=east] at (-0.55, 0) {$\mathit{ok}'(v)$};
    \draw[->, line width=1.1pt, draw=fgblue] (7*0.72, -1.05) node[below, font=\sffamily\scriptsize\bfseries, text=fgblue]
      {last true = (first false) $- 1$} -- (7*0.72, -0.4);
  \end{scope}
\end{tikzpicture}
\caption{Monotone predicates. Top: false up to some point, true afterwards; binary search finds the first true value.
Bottom: the ``maximise'' form, true up to some point; its last true value is one less than its first false value.}
\label{fig:monotone}
\end{figure}

\begin{pseudo}{First true of a monotone predicate}
\begin{algorithmic}[1]
\Require $\mathit{ok}$ monotone on $[lo, hi]$, and $\mathit{ok}(hi)$ true (or $hi$ means ``no answer'')
\While{$lo < hi$} \Comment{invariant: the answer lies in $[lo, hi]$}
  \State $mid \gets lo + \lfloor (hi - lo) / 2 \rfloor$ \Comment{$lo \le mid < hi$}
  \If{$\mathit{ok}(mid)$} \State $hi \gets mid$ \Comment{$mid$ may be the answer}
  \Else \State $lo \gets mid + 1$ \Comment{everything up to $mid$ is false}
  \EndIf
\EndWhile
\State \Return $lo$
\end{algorithmic}
\end{pseudo}

Three details make this loop robust. First, $mid$ is computed as $lo + (hi - lo)/2$, never as $(lo + hi)/2$, which overflows
when both are close to $2^{63}$. Second, because $mid < hi$, both branches strictly shrink the range, so the loop terminates;
the tempting variant \cc{lo = mid} (without $+1$) can leave $lo$ unchanged when $hi = lo + 1$ and loop forever. Third, the
initial $hi$ must be a value where $\mathit{ok}$ is certainly true, or a sentinel one past the range that you check for after
the loop; then the answer is always inside $[lo, hi]$.

For the ``maximise'' form (the last $v$ with $\mathit{ok}(v)$ true, as in Figure~\ref{fig:monotone}, bottom), either search
for the first false value and subtract one, or mirror the loop: round $mid$ \emph{up}, $mid = lo + \lceil (hi - lo)/2\rceil$,
and set $lo \gets mid$ on true and $hi \gets mid - 1$ on false. The next section uses the second form.

\usetemplate{Binary search on the answer}{Ctrl+Alt+T, A}{cpbsearch}

\begin{trick}[Binary search by jumps]
An equivalent formulation, often shorter and immune to off-by-one errors, builds the answer bit by bit: start at a position
where $\mathit{ok}$ is false and take jumps of decreasing powers of two, moving whenever the target of the jump is still
false.
\begin{code}
ll x = lo - 1;  // ok(lo - 1) is false (or lo - 1 is outside the domain)
for (ll step = 1LL << 62; step > 0; step >>= 1)
    if (x + step < hi && !ok(x + step)) x += step;
// x is the last false value, so x + 1 is the first true one
\end{code}
The same pattern descends a Fenwick tree to find a prefix sum in $\Oh{\log n}$ (Chapter~\ref{ch:segtree}) and climbs a tree
with binary lifting (Chapter~\ref{ch:trees}).
\end{trick}

\begin{trick}[Galloping when the range is unknown]
If the answer can be arbitrarily large and $\mathit{ok}$ gets expensive for large values, find an upper bound first by
doubling: try $1, 2, 4, 8, \ldots$ until $\mathit{ok}$ becomes true, then binary search between the last two tries. This
\emph{exponential search} costs $\Oh{\log v}$ evaluations, where $v$ is the answer itself rather than the largest possible
answer, and it never evaluates $\mathit{ok}$ far beyond $v$, which also avoids overflow inside $\mathit{ok}$.
\end{trick}

\subsection{Beyond monotonicity: keep a half that surely contains an answer}

Binary search needs less than a monotone predicate. All it needs is that, after each question, \emph{one half is guaranteed
to contain some answer}. Bisection for roots is the classic example: if a continuous function has $f(lo) < 0 < f(hi)$, then
whichever half has a sign change at its ends contains a root, even though $f$ need not be monotone. The discrete analogue is
just as useful.

\begin{problem}{Searching Local Minimum (Codeforces 1479A)}
An unknown permutation $a_1, \ldots, a_n$ ($n \le 10^5$) can be read one element per query. Find a \emph{local minimum},
an index $i$ with $a_{i-1} > a_i < a_{i+1}$ (where $a_0 = a_{n+1} = \infty$), using at most 100 queries.
\solution Keep a range $[lo, hi]$ with $a_{lo - 1} > a_{lo}$ and $a_{hi} < a_{hi + 1}$: walking downhill from $lo$ into the
range can never leave it, so it contains a local minimum. Initially $[1, n]$ qualifies. Read $a_{mid}$ and $a_{mid+1}$: if
$a_{mid} < a_{mid+1}$, the range $[lo, mid]$ still qualifies; otherwise $[mid + 1, hi]$ does. When $lo = hi$, both
conditions together say that $lo$ is a local minimum (Figure~\ref{fig:localmin}). This takes $2\lceil\log_2 n\rceil \le 34$
queries.
\end{problem}

\begin{figure}[htbp]
\centering
\begin{tikzpicture}[x=0.8cm, y=0.32cm]
  \foreach \i/\v in {1/5, 2/3, 3/4, 4/7, 5/6, 6/2, 7/1, 8/8} {
    \ifnum\i=7 \def\bst{good}\else\ifnum\i=2 \def\bst{warn}\else\def\bst{cool}\fi\fi
    \draw[\bst] (\i - 0.3, 0) rectangle (\i + 0.3, \v);
    \node[font=\sffamily\scriptsize] at (\i, \v + 0.6) {\v};
    \node[lbl] at (\i, -0.7) {\i};
  }
  \foreach \i in {0, 9} {
    \draw[line width=0.6pt, draw=fgline, dashed] (\i - 0.3, 0) -- (\i - 0.3, 9.5) (\i + 0.3, 0) -- (\i + 0.3, 9.5);
    \node[font=\sffamily\scriptsize] at (\i, 10.2) {$\infty$};
    \node[lbl] at (\i, -0.7) {\i};
  }
  \draw[line width=0.6pt, draw=fgink!70] (-0.6, 0) -- (9.6, 0);
  \foreach \s/\a/\b/\h in {1/4/5/8.4, 2/6/7/3.6, 3/7/8/9.6} {
    \draw[->, line width=0.9pt, draw=fgverm] (\a, \h) to[bend left=35] (\b, \h);
    \node[circle, fill=fgverm, text=white, font=\sffamily\tiny\bfseries, inner sep=1pt] at ({(\a + \b)/2}, \h + 1.3) {\s};
  }
  \node[notel, anchor=west, text width=4.6cm] at (10.2, 5) {Step 1: $7 > 6$, go right to $[5, 8]$.\\ Step 2: $2 > 1$, go
    right to $[7, 8]$.\\ Step 3: $1 < 8$, keep $[7, 7]$: index 7 (green) is a local minimum. Index 2 (orange) is another
    one, in the half that was dropped.};
\end{tikzpicture}
\caption{Finding a local minimum by binary search. Each step compares $a_{mid}$ with $a_{mid+1}$ and keeps the half on the
downhill side; that half is guaranteed to contain a local minimum, although not necessarily the global one.}
\label{fig:localmin}
\end{figure}

\begin{code}
map<int, int> memo;
auto ask = [&](int i) -> int {  // a[i], with a[0] = a[n + 1] = infinity
    if (i < 1 || i > n) return 1e9;
    if (!memo.count(i)) {
        cout << "? " << i << endl;  // endl flushes: the judge waits for it
        cin >> memo[i];
    }
    return memo[i];
};
int lo = 1, hi = n;  // a local minimum lies in [lo, hi]
while (lo < hi) {
    int mid = (lo + hi) / 2;
    if (ask(mid) < ask(mid + 1)) hi = mid;
    else lo = mid + 1;
}
cout << "! " << lo << endl;
\end{code}

\begin{tip}[Interactive problems]
In an interactive problem the judge answers your queries while your program runs, so every query must be \emph{flushed}
(\cc{endl} or \cc{cout.flush()}), and the query limit is part of the problem: $\lceil \log_2 n \rceil$ queries allow a plain
binary search, $2\lceil \log_2 n\rceil$ allow comparing neighbours, and a limit like $n/2$ hints at pairing. Cache answers so
that no position is asked twice, and test locally with a small program that plays the judge. The template
\texttt{cpinteractive} (\key{Ctrl+Alt+T, Shift+I}) has the query helpers.
\end{tip}

% ======================================================================================================================
\section[Binary search on the answer]{Binary search on the answer \hfill \silver}\label{sec:bsearch-answer}

Many optimisation problems ask for the smallest (or largest) value with some property, where \emph{checking} whether a given
value works is much easier than \emph{finding} the best value directly. If the property is monotone, binary search on the
answer reduces the optimisation problem to $\Oh{\log(\text{range})}$ checks. This is the idea behind Array Division
(Section~\ref{sec:method}) and Factory Machines (Chapter~\ref{ch:worked}); here is the classic ``maximise the minimum''
version.

\begin{problem}{Aggressive cows}
There are $n$ stalls at given positions on a line. Place $k \ge 2$ cows in distinct stalls so that the minimum distance
between two cows is as large as possible.
\solution Ask instead: \emph{can the cows be placed so that every two are at least $d$ apart?} This is monotone (if a
placement works for $d$, the same placement works for every smaller $d$), and it is easy to check greedily: sort the stalls,
put a cow in the first stall, and put each next cow in the first stall at least $d$ beyond the previous cow. Placing each cow
as early as possible never hurts, by a ``greedy stays ahead'' argument like the one for interval scheduling
(Section~\ref{sec:greedy}). The answer is the largest $d$ for which the check succeeds (Figure~\ref{fig:cows}), found in
$\Oh{n \log n + n \log(\text{range})}$.
\end{problem}

\begin{figure}[htbp]
\centering
\begin{tikzpicture}[x=0.95cm]
  \foreach \yy/\dd/\cows/\res/\col in {0/3/{1,4,8}/{3 cows: \ yes}/fggreen, -1.3/4/{1,8}/{2 cows: \ no}/fgverm} {
    \draw[line width=0.6pt, draw=fgink!60] (0.5, \yy) -- (9.5, \yy);
    \foreach \s in {1, 2, 4, 8, 9} {
      \draw[line width=0.6pt, draw=fgink!70] (\s, \yy - 0.1) -- (\s, \yy + 0.1);
      \node[lbl, below=1pt] at (\s, \yy - 0.08) {\s};
    }
    \foreach \cw in \cows \node[circle, fill=\col, inner sep=0pt, minimum size=3.2mm] at (\cw, \yy + 0.32) {};
    \node[blbl, anchor=east] at (0.3, \yy) {$d = \dd$};
    \node[font=\sffamily\scriptsize\bfseries, text=\col, anchor=west] at (9.8, \yy) {\res};
  }
\end{tikzpicture}
\caption{The greedy check for stalls $1, 2, 4, 8, 9$ and $k = 3$ cows. With $d = 3$ the cows go to 1, 4 and 8; with
$d = 4$ only two fit (1 and 8). The answer is 3, the largest $d$ that passes.}
\label{fig:cows}
\end{figure}

\begin{code}
void solve() {
    int n, k;
    cin >> n >> k;
    vi x(n);
    cin >> x;
    sort(all(x));
    // can k cows be placed with every two at least d apart? (greedy)
    auto ok = [&](ll d) {
        int cows = 1;
        ll last = x[0];
        rep(i, 1, n) if (x[i] - last >= d) cows++, last = x[i];
        return cows >= k;
    };
    ll lo = 0, hi = x[n - 1] - x[0];  // ok(lo) holds; find the last d with ok
    while (lo < hi) {
        ll mid = lo + (hi - lo + 1) / 2;  // round up, or lo = mid never moves
        if (ok(mid)) lo = mid;
        else hi = mid - 1;
    }
    print(lo);
}
\end{code}

\begin{recognize}
``Minimise the maximum'', ``maximise the minimum'', ``the smallest time, size or capacity such that\ldots'', or any answer
where checking a candidate is easier than constructing the best one. The check is often a greedy simulation, and the
monotonicity argument is usually one sentence: ``a solution for the stricter requirement is also a solution for the looser
one''.
\end{recognize}

\subsection{The $k$-th smallest element of an implicit set}

A second family of problems asks for the $k$-th smallest element of a set that is far too large to build: all pairwise sums,
all products in a table, all subarray sums. Binary search the \emph{value}: the $k$-th smallest is the smallest $x$ such that
at least $k$ elements are $\le x$. Counting the elements $\le x$ is usually much easier than listing them.

\begin{problem}{Multiplication Table (CSES)}
Find the median of the $n \times n$ multiplication table (the entries $i \cdot j$ for $1 \le i, j \le n$), where $n \le 10^6$
is odd.
\solution The median is the $k$-th smallest entry for $k = (n^2 + 1)/2$. Row $i$ contains the multiples $i, 2i, \ldots, ni$,
and exactly $\min(n, \lfloor x / i\rfloor)$ of them are at most $x$, so counting the entries $\le x$ takes $\Oh{n}$
(Figure~\ref{fig:multtable}). Binary search the smallest $x$ whose count reaches $k$: $\Oh{n \log n^2}$ in total, 0.07\,s for
$n = 10^6$. The answer is always an entry of the table, because the count only increases at entries.
\end{problem}

\begin{figure}[htbp]
\centering
\begin{tikzpicture}
  \foreach \i in {1,...,5} \foreach \j in {1,...,5} {
    \pgfmathtruncatemacro{\pv}{\i*\j}
    \ifnum\pv<9 \node[cellc, cool, minimum width=7.5mm, minimum height=6.5mm] at (\j*0.75, -\i*0.65) {\pv};
    \else \node[cellc, minimum width=7.5mm, minimum height=6.5mm] at (\j*0.75, -\i*0.65) {\pv}; \fi
  }
  \foreach \i/\c in {1/5, 2/4, 3/2, 4/2, 5/1} \node[font=\sffamily\small\bfseries, text=fgblue] at (4.7, -\i*0.65) {\c};
  \node[lbl] at (4.7, -0.1) {$\le 8$};
  \node[notel, anchor=west, text width=6.2cm] at (5.5, -1.95) {For $x = 8$: $5 + 4 + 2 + 2 + 1 = 14 \ge 13$ entries are at
    most 8; for $x = 7$ only 12. So the median, the 13th smallest entry, is 8.};
\end{tikzpicture}
\caption{The $5 \times 5$ multiplication table. The entries $\le 8$ (blue) form a staircase; row $i$ contributes
$\min(n, \lfloor 8/i \rfloor)$ of them (right).}
\label{fig:multtable}
\end{figure}

\begin{code}
void solve() {
    ll n;
    cin >> n;
    ll k = (n * n + 1) / 2;  // the median is the k-th smallest entry
    // how many entries i * j (1 <= i, j <= n) are at most x?
    auto count = [&](ll x) {
        ll c = 0;
        for (ll i = 1; i <= n; i++) c += min(n, x / i);
        return c;
    };
    ll lo = 1, hi = n * n;  // the first x with count(x) >= k
    while (lo < hi) {
        ll mid = lo + (hi - lo) / 2;
        if (count(mid) >= k) hi = mid;
        else lo = mid + 1;
    }
    print(lo);
}
\end{code}

The same pattern finds the $k$-th smallest pairwise sum (count pairs with sum $\le x$ by two pointers), the $k$-th smallest
subarray sum of a positive array (two pointers again), and the $k$-th smallest distance between points on a line. When the
count itself needs a data structure, the whole search is sometimes replaced by a walk down that structure, as in the
order-statistics tree of Chapter~\ref{ch:advds}.

% ======================================================================================================================
\section[Real numbers and fractional programming]{Real numbers and fractional programming \hfill \gold}\label{sec:bsearch-real}

When the answer is a real number, the loop cannot run until $lo = hi$. The robust approach is to run a \emph{fixed number of
iterations}: each halves the interval, so 100 iterations shrink any interval of \cc{double}s far below their precision.
Comparing \cc{hi - lo} with an epsilon is fragile instead: with values around $10^{9}$ and \cc{eps = 1e-9} the loop never
ends, because consecutive doubles near $10^9$ are about $10^{-7}$ apart and the interval cannot shrink further.

\usetemplate{Binary search on a real answer}{Ctrl+Alt+T, Shift+F}{cpbsreal}

\subsection{Maximising a ratio}

A particularly important use of real-valued binary search is \defn{fractional programming}: maximising a ratio
$\sum a_i / \sum b_i$ (with all $b_i > 0$) over some family of choices, such as subsets, paths, cycles or subarrays. Ratios do
not decompose, so DP and greedy algorithms cannot optimise them directly. But for any $\lambda$,
\[
  \frac{\sum a_i}{\sum b_i} \ge \lambda \iff \sum (a_i - \lambda b_i) \ge 0 ,
\]
so ``can the ratio reach $\lambda$?'' is the question ``is the maximum of $\sum (a_i - \lambda b_i)$ non-negative?'', an
ordinary optimisation problem with weights $a_i - \lambda b_i$. The answer is monotone in $\lambda$, so binary search on
$\lambda$ finds the best ratio.

\begin{problem}{Maximum average subarray}
Given $n \le 10^5$ numbers and $k$, find the maximum average of a subarray of length at least $k$.
\solution The average is the ratio $\sum a_i / \sum 1$. For a candidate $\lambda$, subtract $\lambda$ from every element and
ask whether some subarray of length $\ge k$ has a non-negative sum. With prefix sums $p$ of the shifted values, that is
whether $p_r - \min_{l \le r - k} p_l \ge 0$ for some $r$, one linear scan. A hundred iterations of binary search then take
$\Oh{100\, n}$.
\end{problem}

\begin{code}
void solve() {
    int n, k;
    cin >> n >> k;
    vector<double> a(n);
    cin >> a;
    // is there a subarray of length >= k with average >= lam? Exactly when
    // some such subarray has sum of (a[i] - lam) >= 0
    auto ok = [&](double lam) {
        vector<double> p(n + 1, 0);
        rep(i, n) p[i + 1] = p[i] + a[i] - lam;
        double lowest = 1e18;  // the smallest p[l] with l <= r - k
        rep(r, k, n + 1) {
            chmin(lowest, p[r - k]);
            if (p[r] - lowest >= 0) return true;
        }
        return false;
    };
    double lo = *min_element(all(a)), hi = *max_element(all(a));
    rep(it, 100) {
        double mid = (lo + hi) / 2;
        if (ok(mid)) lo = mid;
        else hi = mid;
    }
    cout << fixed << setprecision(6) << lo << '\n';
}
\end{code}

\paragraph{Dinkelbach's method.} Write $g(\lambda) = \max \sum (a_i - \lambda b_i)$ over all choices. Each choice contributes
the line $A - \lambda B$ (with $A = \sum a_i$, $B = \sum b_i$), so $g$ is the upper envelope of lines with negative slopes:
decreasing and convex, with its root exactly at the best ratio $\lambda^\ast$ (Figure~\ref{fig:dinkelbach}). Instead of
bisecting, \defn{Dinkelbach's method} jumps: solve the inner problem at the current $\lambda$, take the optimal choice's own
ratio $A/B$ as the next $\lambda$, and repeat until $g(\lambda) = 0$. This is Newton's method on $g$, and it typically
converges in a handful of iterations, and exactly (with no precision issues) when the inner problem can be solved with
fractions or integers.

\begin{figure}[htbp]
\centering
\begin{tikzpicture}
  \begin{axis}[cpplot, width=10.5cm, height=6cm, xmin=0, xmax=5.2, ymin=-6, ymax=11, xlabel={$\lambda$},
      ylabel={$A - \lambda B$}, xtick={0,1,2,2.5,3,4,5}, xticklabels={0,1,2,2.5,3,4,5}, ytick={-5,0,5,10},
      legend pos=north east, extra y ticks={0}, extra y tick style={grid style={fgink!60, line width=0.6pt}}]
    \addplot[fgline, line width=0.9pt, domain=0:5.2, samples=2] {6 - 3*x};
    \addlegendentry{$A{=}6, B{=}3$}
    \addplot[fgorange!90!black, line width=0.9pt, domain=0:5.2, samples=2] {10 - 4*x};
    \addlegendentry{$A{=}10, B{=}4$}
    \addplot[fggreen, line width=0.9pt, domain=0:5.2, samples=2] {4 - x};
    \addlegendentry{$A{=}4, B{=}1$}
    \addplot[fgblue, line width=2.4pt, opacity=0.35, domain=0:2, samples=2, forget plot] {10 - 4*x};
    \addplot[fgblue, line width=2.4pt, opacity=0.35, domain=2:5.2, samples=2, forget plot] {4 - x};
    \addplot[only marks, mark=*, mark size=2.4pt, fgverm, forget plot] coordinates {(2.5, 0) (4, 0)};
    \draw[->, line width=0.9pt, draw=fgverm] (axis cs:2.5, 0.2) -- (axis cs:2.5, 1.45);
    \node[font=\sffamily\scriptsize, text=fgverm, anchor=north] at (axis cs:2.5, -0.35) {$\lambda_1$};
    \node[font=\sffamily\scriptsize, text=fgverm, anchor=north] at (axis cs:4, -0.35) {$\lambda^\ast = \lambda_2$};
  \end{axis}
\end{tikzpicture}
\caption{Fractional programming with three choices. The best ratio $\lambda^\ast = 4$ is the root of the upper envelope $g$
(thick blue). Dinkelbach's method starts at $\lambda_0 = 0$, where the best line is $10 - 4\lambda$; its root gives
$\lambda_1 = 10/4$. At $\lambda_1$ the best line is $4 - \lambda$, whose root $\lambda_2 = 4$ is the answer.}
\label{fig:dinkelbach}
\end{figure}

Fractional programming appears in disguise surprisingly often: the minimum mean cycle (weights $w_e - \lambda$, check for a
negative cycle), the optimal ratio spanning tree (MST with weights $a_e - \lambda b_e$), the densest subgraph (maximum
closure with weights shifted by $\lambda$), and ``best average'' problems of every kind.

% ======================================================================================================================
\section[Ternary search]{Ternary search on unimodal functions \hfill \gold}\label{sec:ternary}

A function is \defn{unimodal} if it strictly increases up to its maximum and strictly decreases afterwards (or the reverse,
for a minimum). To find the peak, evaluate $f$ at two interior points $m_1 < m_2$. If $f(m_1) < f(m_2)$, the peak cannot be to
the left of $m_1$: there the function is still increasing, and a peak to the left would make $f$ decrease from $m_1$ to
$m_2$. Otherwise the peak cannot be to the right of $m_2$. Either way a third of the interval is discarded
(Figure~\ref{fig:ternary}), and after $k$ steps the interval has shrunk by $(2/3)^k$: 100 steps are plenty for doubles.

\begin{figure}[htbp]
\centering
\begin{tikzpicture}
  \begin{axis}[cpplot, width=10cm, height=5.2cm, xmin=0, xmax=10, ymin=0, ymax=11.5, xtick={0,10}, ytick=\empty,
      xlabel={$x$}, ylabel={$f(x)$}, grid=none, clip=false]
    \fill[fgline!25] (axis cs:0, 0) rectangle (axis cs:3.33, 11.5);
    \addplot[fgblue, line width=1.6pt, domain=0.86:10, samples=100, mark=none] {10 - 0.35*(x-6.2)^2};
    \addplot[only marks, mark=*, mark size=2.6pt, fgverm] coordinates {(3.33, 7.12) (6.67, 9.92)};
    \draw[dashed, line width=0.7pt, draw=fgverm] (axis cs:3.33, 0) -- (axis cs:3.33, 7.12);
    \draw[dashed, line width=0.7pt, draw=fgverm] (axis cs:6.67, 0) -- (axis cs:6.67, 9.92);
    \node[font=\sffamily\scriptsize, text=fgverm, anchor=north] at (axis cs:3.33, -0.2) {$m_1$};
    \node[font=\sffamily\scriptsize, text=fgverm, anchor=north] at (axis cs:6.67, -0.2) {$m_2$};
    \node[font=\sffamily\scriptsize, text=fggray] at (axis cs:1.65, 10.5) {discarded};
    \node[font=\sffamily\scriptsize, text=fgink, anchor=west] at (axis cs:7.1, 10.6) {$f(m_1) < f(m_2)$};
  \end{axis}
\end{tikzpicture}
\caption{One step of ternary search for a maximum. Since $f(m_1) < f(m_2)$, the maximum cannot lie in $[0, m_1]$ (grey),
which is discarded.}
\label{fig:ternary}
\end{figure}

\paragraph{Integers: binary search on the slope.} For a function on integers, ternary search needs care near the end. A
cleaner method uses the fact that a unimodal function's \emph{difference} $f(x+1) - f(x)$ changes sign once: binary search the
first $x$ with $f(x) \ge f(x + 1)$ (for a maximum). This is an ordinary monotone predicate and costs $2\log(\text{range})$
evaluations. The template contains both versions.

\usetemplate{Ternary search (unimodal max)}{Ctrl+Alt+T, Q}{cpternary}

\paragraph{Convexity is the real condition.} In practice the functions you ternary search are usually \emph{convex}:
maxima of linear functions, sums of absolute values, sums of convex costs. A convex function is unimodal (or constant on its
minimum), and sums and maxima of convex functions are convex, so convexity is easy to establish from the structure of the
cost even when the function has no closed form.

\begin{problem}{Weakness and Poorness (Codeforces 578C)}
Given $n \le 2\cdot10^5$ integers $a_i$ with $|a_i| \le 10^4$, find a real $x$ minimising the \emph{weakness} of the sequence
$a_1 - x, \ldots, a_n - x$: the largest absolute value of the sum of a segment.
\solution For a fixed segment, its sum $S - Lx$ is linear in $x$, so its absolute value is convex; the weakness is the
maximum of these convex functions, hence convex in $x$. Ternary search on $x$; for a given $x$, the largest segment sum and
the smallest segment sum come from two runs of Kadane's algorithm (Chapter~\ref{ch:worked}), and the weakness is the larger
of the maximum and minus the minimum.
\end{problem}

\begin{code}
void solve() {
    int n;
    cin >> n;
    vector<double> a(n);
    cin >> a;
    // the largest |sum| of a segment of a[i] - x: the larger of the best
    // sum and minus the worst sum (Kadane); a convex function of x
    auto weakness = [&](double x) {
        double best = 0, worst = 0, hi = 0, lo = 0;
        for (double v : a) {
            hi = max(0.0, hi) + v - x;
            lo = min(0.0, lo) + v - x;
            best = max(best, hi);
            worst = min(worst, lo);
        }
        return max(best, -worst);
    };
    double l = -1e4, r = 1e4;
    rep(it, 100) {  // ternary search: each step keeps 2/3 of the range
        double m1 = l + (r - l) / 3, m2 = r - (r - l) / 3;
        if (weakness(m1) < weakness(m2)) r = m2;
        else l = m1;
    }
    cout << fixed << setprecision(9) << weakness((l + r) / 2) << '\n';
}
\end{code}

\begin{pitfall}[Plateaus break ternary search]
If $f$ is constant on some interval other than its optimum, $f(m_1) = f(m_2)$ says nothing about where the optimum is, and
ternary search may discard it. Ternary search is safe for strictly unimodal functions and for convex functions (whose only
possible plateau is the set of minima). For functions of two variables, nested ternary search works when the function is
convex in both jointly, but is only a heuristic otherwise.
\end{pitfall}

\begin{trick}[Golden-section search]
Ternary search evaluates $f$ twice per step. If the two interior points divide the interval in the golden ratio,
$m_1 = r - (r - l)/\varphi$ and $m_2 = l + (r - l)/\varphi$ with $\varphi = (1 + \sqrt 5)/2$, then after discarding one
side, the surviving interior point is exactly one of the two points needed in the next step. So each step costs \emph{one}
evaluation and shrinks the interval by $1/\varphi \approx 0.618$: about 1.4 times fewer evaluations than ternary search for the
same precision, which matters when each evaluation is a full $\Oh{n}$ pass.
\end{trick}

% ======================================================================================================================
\section[Parallel binary search]{Parallel binary search \hfill \ioi}\label{sec:parallel-bs}

Suppose there are $q$ queries, each of which could be answered by its own binary search over the same sequence of $m$ events:
``after how many events does query $i$ become true?''. Each check replays events from the beginning, so separate searches cost
$\Oh{q \cdot m \log m}$ event operations. \defn{Parallel binary search} runs all $q$ searches together, in rounds. In each
round, every unfinished query has a current range and a middle; sort the queries by their middles, replay the events
\emph{once} from the start, and check each query at the moment the replay reaches its middle. After $\lceil \log_2 (m + 2)
\rceil$ rounds every range has shrunk to one value, for a total of $\Oh{(m + q)\log m}$ event operations plus sorting.

\begin{problem}{First moment of connectivity}
$m$ edges are added one by one to a graph with $n$ vertices. For each of $q$ queries $(a, b)$, find the first time after which
$a$ and $b$ are connected (0 if $a = b$, $-1$ if never). Limits: $n, m, q \le 2\cdot10^5$.
\solution Connectivity only grows as edges are added, so each query's answer is a monotone boundary. In each round, sort the
queries by their middles, rebuild a DSU from scratch and add the edges in order, checking each query when its middle is
reached (Figure~\ref{fig:pbs}). Eighteen rounds of $m$ unions and $q$ checks answer everything in 0.4\,s.
\end{problem}

\begin{figure}[htbp]
\centering
\small
\begin{tabular}{@{}l c c c c@{}}
  \toprule
  \thd{Query (answer)} & \thd{Round 1} & \thd{Round 2} & \thd{Round 3} & \thd{Round 4} \\
  \midrule
  $q_1$ (3) & $[0, 8]$, mid 4: \textcolor{fggreen}{yes} & $[0, 4]$, mid 2: \textcolor{fgverm}{no} &
    $[3, 4]$, mid 3: \textcolor{fggreen}{yes} & $[3, 3]$ done \\
  $q_2$ (6) & $[0, 8]$, mid 4: \textcolor{fgverm}{no} & $[5, 8]$, mid 6: \textcolor{fggreen}{yes} &
    $[5, 6]$, mid 5: \textcolor{fgverm}{no} & $[6, 6]$ done \\
  $q_3$ (1) & $[0, 8]$, mid 4: \textcolor{fggreen}{yes} & $[0, 4]$, mid 2: \textcolor{fggreen}{yes} &
    $[0, 2]$, mid 1: \textcolor{fggreen}{yes} & $[0, 1]$, mid 0: \textcolor{fgverm}{no} \\
  $q_4$ (never) & $[0, 8]$, mid 4: \textcolor{fgverm}{no} & $[5, 8]$, mid 6: \textcolor{fgverm}{no} &
    $[7, 8]$, mid 7: \textcolor{fgverm}{no} & $[8, 8]$ done \\
  \midrule
  \thd{Replays} & edges 1--4 & edges 1--6 & edges 1--7 & none \\
  \bottomrule
\end{tabular}
\caption{Parallel binary search with $m = 7$ events, so every range starts as $[0, 8]$ (8 means ``never''). In each round the
queries are checked in order of their middles during a single replay of the events; ``yes'' means connected after
\emph{mid} events. After round 4 every range is a single value; for $q_3$ it is $[1, 1]$.}
\label{fig:pbs}
\end{figure}

\begin{code}
void solve() {
    int n, m, q;
    cin >> n >> m >> q;
    vector<pii> e(m), qs(q);
    for (auto& [u, v] : e) cin >> u >> v;
    for (auto& [a, b] : qs) cin >> a >> b;
    vi lo(q, 0), hi(q, m + 1);  // answer in [lo, hi]; m + 1 means never
    while (true) {
        vector<pii> byMid;  // (mid, query) for the unfinished queries
        rep(i, q) if (lo[i] < hi[i]) byMid.push_back({(lo[i] + hi[i]) / 2, i});
        if (byMid.empty()) break;
        sort(all(byMid));
        DSU d(n + 1);  // one replay of the edges serves every query
        int added = 0;
        for (auto [mid, i] : byMid) {
            while (added < mid) {  // add edges 1..mid
                d.unite(e[added].first, e[added].second);
                added++;
            }
            if (d.same(qs[i].first, qs[i].second)) hi[i] = mid;
            else lo[i] = mid + 1;
        }
    }
    rep(i, q) print(lo[i] > m ? -1 : lo[i]);
}
\end{code}

The technique needs only two things: each query's answer is monotone in the number of events, and the events can be replayed
from scratch cheaply. It shines when the events are hard to \emph{undo} (DSU unions, which cannot be rolled back without extra
work) or when the per-query check needs a data structure that all queries can share. The classic example is Meteors
(POI 2011): each of $k$ meteor showers adds rain to a range of sectors, and each country wants to know after which shower its
sectors have collected enough; each round replays the showers with a Fenwick tree for range additions, for
$\Oh{(n + k)\log k \log n}$ in total.

\begin{summary}[Chapter summary]
\begin{itemize}
  \item Write binary search from an invariant: the answer stays in $[lo, hi]$, $mid = lo + (hi - lo)/2$, and one branch
        moves $lo$ past $mid$.
  \item \texttt{lower\_bound} is the first $\ge x$, \texttt{upper\_bound} the first $> x$; on a \texttt{set} use the member
        functions.
  \item Binary search on the answer turns optimisation into checking; the $k$-th smallest of an implicit set is the first
        value with at least $k$ elements below it.
  \item For real answers run a fixed number of iterations; fractional programming turns ratios into sums with weights
        $a_i - \lambda b_i$, and Dinkelbach's method converges in a few steps.
  \item Ternary search optimises unimodal (in practice, convex) functions; on integers, binary search the slope instead.
  \item Parallel binary search answers many monotone queries with one replay of the events per round.
\end{itemize}
\end{summary}

\begin{practice}
\ppgroup{Binary search}
\begin{pplist}
  \item Concert Tickets
  \item Sum of Two Values
  \item Factory Machines
  \item Array Division
  \item Multiplication Table
  \item Increasing Subsequence
\end{pplist}
\end{practice}
